% chapters/ch4-symbols.tex -- Chapter 4 Symbols list (PDF pages 531-536)
% Corrections applied (corrections/ch4.md): item 6 from Mandell's ``Changes to text of Chapter 4'' (June 1994,
% PDF 702): six entries added in the list's order, each vector next to its scalar (\vec{A}_e after A, B;
% \vec{E}_o after \vec{E}; \vec{F}(t) after F(t); P_a, P_e after M_y; \vec{c}_eff after \vec{c}), one note
% on the first of them; the meaning of F(t) changed, with a note quoting the 1973 entry.
% An \ednote in a p{} cell of a longtable loses its text (STYLE.md section 9): the F(t) row is a one-line
% \multicolumn{1}{l@{}} cell; the \vec{A}_e row is too long for one line and carries its note as \edcap.
\unnumberedsection{Symbols}\label{ch4:sec:symbols}
\begin{longtable}{@{}l p{0.78\linewidth}@{}}
\toprule \textbf{Symbol} & \textbf{Meaning} \\ \midrule \endfirsthead
\toprule \textbf{Symbol} & \textbf{Meaning} \\ \midrule \endhead
\bottomrule \endlastfoot
$A, B$ & dummy variables used in trigonometric identities \\
$\vec{A}_e$ & nozzle exit plane area written as a vector whose direction is forward along the vehicle centerline \edcap{The symbols $\vec{A}_e$, $\vec{E}_o$, $\vec{F}(t)$, $P_a$, $P_e$ and $\vec{c}_{\mathrm{eff}}$ are added per Mandell's June 1994 supplement (``Changes to text of Chapter~4'', p.~3); they do not appear in the 1973 Symbols list. They are introduced in the corrected text after equation~\eqref{ch4:eq:7}.} \\
$A_f$ & amplitude of sinusoidal forcing \\
$A_r$ & reference area \\
$A_o$ & value of $A_f$ at an airspeed of 1 meter/second \\
$A_1, A_2$ & coefficients used in Riccati solution \\
$\CD$ & coefficient of drag \\
$C_1$ & corrective moment coefficient \\
$C_2$ & damping moment coefficient \\
$D$ & drag \\
$\vec{E}$ & vector sum of all externally applied forces \\
$\vec{E}_o$ & vector sum of all externally-applied forces other than pressure component of thrust \\
$F$ & average thrust \\
$\vec{F}$ & force \\
$F(t)$ & \multicolumn{1}{l@{}}{magnitude of thrust as a function of time\ednote{Changed per Mandell's June 1994 supplement (``Changes to text of Chapter~4'', p.~3), which adds the thrust vector $\vec{F}(t)$ as a separate entry. The 1973 entry read ``thrust as a function of time''.}} \\
$\vec{F}(t)$ & thrust as a function of time \\
$F_m$ & maximum thrust \\
$F_n$ & average thrust of $n$th stage motor \\
$F_p$ & component of thrust perpendicular to trajectory \\
$F_s$ & sustainer thrust \\
$F_t$ & component of thrust tangent to trajectory \\
$F_2$ & average thrust of second stage motor \\
$H$ & function of thrust, mass, and drag parameter used in writing Riccati solution \\
$H_x$ & strength of yaw impulse \\
$H_y$ & strength of pitch impulse \\
$I_L$ & longitudinal moment of inertia \\
$I_R$ & radial moment of inertia \\
$\Isp$ & specific impulse \\
$I_t$ & total impulse \\
$M_x$ & strength of step moment in yaw \\
$M_y$ & strength of step moment in pitch \\
$P_a$ & ambient pressure \\
$P_e$ & pressure of exhaust gas at nozzle exit plane \\
$\vec{a}$ & acceleration \\
$c$ & magnitude of exhaust velocity \\
$\vec{c}$ & exhaust velocity \\
$\vec{c}_{\mathrm{eff}}$ & effective exhaust velocity, \emph{also} called equivalent exhaust velocity \\
$d(\ )/dt$ & derivative of ( ) with respect to time \\
$e$ & base of the Napierian logarithm system, numerically equal to approximately 2.718 \\
$f(\alpha)$ & function of angle of attack \\
$f_x(t)$ & yaw forcing function \\
$f_y(t)$ & pitch forcing function \\
$g$ & acceleration of gravity \\
$k$ & parameter of drag at zero angle of attack \\
$k_n$ & drag parameter of $n$th stage carrying all subsequent stages atop it \\
$k_2$ & drag parameter of second stage \\
$m$ & mass; \emph{also} average mass \\
$m(t)$ & mass as a function of time \\
$m_b$ & burnout mass \\
$m_f$ & mass of propellant charge \\
$m_n$ & mass of $n$th stage carrying all subsequent stages atop it \\
$m_o$ & initial mass \\
$m_1$ & average mass of first stage carrying all subsequent stages atop it \\
$m_2$ & average mass of second stage \\
$\mdot$ & mass flow rate (alternate notation) \\
$dm_e$ & differential of expelled mass \\
$dm_e/dt$ & mass flow rate \\
$\vec{p}$ & momentum \\
$d\vec{p}$ & differential change in momentum \\
$d\vec{p}_E$ & differential change in momentum due to the action of externally applied forces \\
$d\vec{p}_e$ & differential quantity of momentum added to the exhaust stream by the rocket motor \\
$t$ & time \\
$t_b$ & burning time \\
$t_c$ & coasting time \\
$t_m$ & time at which maximum thrust occurs \\
$t_n$ & burning time of $n$th stage rocket motor \\
$t_o$ & time at which in-flight disturbance begins to act \\
$t_s$ & time at which sustainer thrust begins \\
$t_1$ & first stage burning time \\
$t_2$ & second stage burning time \\
$dt$ & differential of time \\
$\Delta t$ & time increment \\
$u$ & variable of transformation used in Riccati analysis \\
$v$ & magnitude of velocity \\
$\vec{v}$ & velocity \\
$v_b$ & burnout velocity \\
$v_n$ & burnout velocity of $n$th stage \\
$v_x$ & horizontal component of velocity \\
$v_y$ & vertical component of velocity \\
$v_1$ & burnout velocity of first stage \\
$v_2$ & burnout velocity of second stage \\
$dv$ & differential change in velocity \\
$\Delta v$ & change in velocity \\
$\Delta v_a$ & actual velocity increment \\
$\Delta v_t$ & drag-free velocity increment \\
$x$ & horizontal coordinate, range \\
$x_b$ & range at burnout \\
$\Delta x$ & increment of range \\
$\dot{x}$ & horizontal component of velocity (alternate notation) \\
$\Delta\dot{x}_a$ & actual increment in horizontal velocity component \\
$\Delta\dot{x}_t$ & drag-free increment in horizontal velocity component \\
$y$ & vertical coordinate, altitude \\
$y_b$ & burnout altitude \\
$y_c$ & coasted altitude increment \\
$y_n$ & altitude increment gained during $n$th stage burn \\
$y_1$ & first stage burnout altitude \\
$y_2$ & altitude increment gained during second stage burn \\
$\Delta y$ & increment in altitude \\
$\dot{y}$ & vertical velocity component (alternate notation) \\
$\Delta\dot{y}_a$ & actual vertical velocity increment \\
$\Delta\dot{y}_t$ & drag-free vertical velocity increment \\
$\alpha$ & angle of attack \\
$\alpha(t)$ & angle of attack as a function of time \\
$\alpha_x$ & yaw angle \\
$\alpha_x(t)$ & yaw angle as a function of time \\
$\alpha_{xo}$ & yaw angle at time $t_o$ \\
$\Delta\alpha_x$ & increment in yaw angle \\
$\alpha_y$ & pitch angle \\
$\alpha_y(t)$ & pitch angle as a function of time \\
$\alpha_{yo}$ & pitch angle at time $t_o$ \\
$\Delta\alpha_y$ & increment in pitch angle \\
$\gamma$ & average rate of mass expulsion \\
$\epsilon$ & parameter of drag due to angle of attack \\
$\zeta$ & damping ratio \\
$\theta$ & angle between tangent to trajectory and local vertical \\
$\theta_o$ & angle between direction of launch and local vertical \\
$\rho$ & density of the atmosphere \\
$\omega$ & angular frequency in radians per second \\
$\omega_{\mathrm{cres}}$ & coupled resonant angular frequency \\
$\omega_f$ & angular frequency of sinusoidal forcing \\
$\omega_n$ & natural frequency \\
$\omega_x$ & yaw rate \\
$\omega_{xo}$ & yaw rate at time $t_o$ \\
$\Delta\omega_x$ & increment in yaw rate \\
$\omega_y$ & pitch rate \\
$\omega_{yo}$ & pitch rate at time $t_o$ \\
$\Delta\omega_y$ & increment in pitch rate \\
$\omega_z$ & roll rate \\
$\omega_{zo}$ & roll rate at an airspeed of 1 meter/second \\
$\omega_o$ & value of $\omega_f$ at an airspeed of 1 meter/second \\
$(\dot{\ })$ & derivative of ( ) with respect to time (alternate notation) \\
$(\ddot{\ })$ & second derivative of ( ) with respect to time (alternate notation) \\
$\int [\ ]\, d(\ )$ & integral of [ ] with respect to ( ) \\
\end{longtable}
\addtocounter{table}{-1}% a caption-less longtable still steps the table counter (STYLE.md section 8)
